\chapter{DỮ LIỆU VÀ PHƯƠNG PHÁP}
\label{chap:methods}

\section{Bộ dữ liệu Flickr8k}

Flickr8k \autocite{hodosh2013framing} gồm 8.000 ảnh đời thường, mỗi ảnh được 5
người mô tả độc lập bằng tiếng Anh. Đồ án dùng bản phân chia Karpathy
\autocite{karpathy2015deep} có sẵn trên Hugging Face (\texttt{jxie/flickr8k}):
6.000 ảnh huấn luyện, 1.000 kiểm định, 1.000 kiểm tra. Mỗi cặp (ảnh, 1 caption)
là một mẫu huấn luyện, cho $6.000\times 5 = 30.000$ mẫu.

\section{Tiền xử lý}

\subsection{Văn bản}
Caption được chuyển chữ thường, loại ký tự không phải chữ/số, tách theo khoảng
trắng. Từ điển mức từ giữ các từ xuất hiện $\geq 5$ lần trong tập huấn luyện,
cộng 4 token đặc biệt \texttt{<pad>}, \texttt{<bos>}, \texttt{<eos>},
\texttt{<unk>}. Đo trên 30.000 caption của tập huấn luyện: từ điển thu được
\textbf{2.539} từ (min\_freq = 5); \textbf{98,4\%} caption có độ dài $\leq 20$
từ (ngưỡng cắt \texttt{max\_words}); sau khi cắt, tỷ lệ token \texttt{<unk>}
trên toàn bộ token huấn luyện là \textbf{2,48\%} --- đủ thấp để ngưỡng cắt và
tần suất tối thiểu không làm mất nhiều thông tin ngữ nghĩa. Chọn mức từ thay vì
BPE vì kho từ vựng Flickr8k nhỏ và mỗi token là một từ trọn vẹn, thuận tiện khi
đọc bản đồ nhiệt attention.

\subsection{Ảnh: trích đặc trưng một lần}
Toàn bộ ảnh được đưa qua ResNet-50 \autocite{he2016resnet} huấn luyện sẵn trên
ImageNet, đóng băng, cắt sau khối \texttt{layer4}: mỗi ảnh cho bản đồ đặc trưng
$7\times 7\times 2048$, dàn phẳng thành 49 vùng $\times$ 2048 chiều, lưu
\texttt{float16}. Quyết định \emph{tính trước một lần} này khiến mọi thí nghiệm
sau đó chỉ huấn luyện decoder trên tensor có sẵn --- mỗi epoch tính bằng giây ---
đổi lại không dùng được augmentation ảnh (nêu ở mục hạn chế).

\section{Kiến trúc hai decoder}

Cả hai decoder nhận cùng đầu vào: 49 vùng ảnh được chiếu tuyến tính
$2048 \to d_{\text{model}} = 512$ (kèm LayerNorm + dropout), và chuỗi từ đã sinh.
Cả hai đều tie trọng số embedding với tầng đầu ra. Khác biệt \emph{duy nhất} là
cơ chế sinh chuỗi --- bảo đảm so sánh công bằng.

\subsection{LSTM với attention Bahdanau}
Một tầng \texttt{LSTMCell} ẩn 512 chiều, khởi tạo trạng thái từ trung bình 49
vùng. Tại mỗi bước, attention cộng tính (công thức~\ref{eq:bahdanau}) dùng trạng
thái ẩn bước trước để tính vector ngữ cảnh; đầu vào của cell là ghép nối
[embedding từ trước; ngữ cảnh]. Biến thể ablation tắt attention thay ngữ cảnh
bằng trung bình cố định của 49 vùng.

\subsection{Transformer decoder tự cài đặt}
Ba khối pre-norm, mỗi khối: masked self-attention (8 đầu) $\to$ cross-attention
vào 49 vùng ảnh (8 đầu) $\to$ FFN 2048. Positional embedding học được; LayerNorm
cuối trước tầng đầu ra. Lý do tự cài từng khối thay vì dùng
\texttt{nn.TransformerDecoder}: lớp đóng hộp không trả ra trọng số
cross-attention, trong khi bản đồ nhiệt ``từ nào nhìn vùng nào'' là phân tích
định tính trung tâm của đồ án.

\begin{table}[H]
	\centering
	\caption{Cấu hình hai decoder (số tham số in trực tiếp từ
	\texttt{model.count\_parameters()} khi khởi tạo, vocab 2.539 từ).}
	\label{tab:architecture}
	\begin{tabular}{lcc}
		\hline
		 & \textbf{LSTM+Bahdanau} & \textbf{Transformer} \\
		\hline
		Số tầng & 1 & 3 \\
		$d_{\text{model}}$ & 512 & 512 \\
		Attention & cộng tính, 1 đầu & tích vô hướng, 8 đầu \\
		FFN & --- & 2048 \\
		Dropout & 0,3 & 0,3 \\
		Tie embedding--output & có & có \\
		Số tham số & 6.551.041 & 14.974.464 \\
		\hline
	\end{tabular}
\end{table}

Biến thể ablation \emph{LSTM không attention} (mục hạn chế attention bằng
mean-pool cố định, xem \ref{tab:main-results}) có \textbf{6.025.216} tham số ---
ít hơn bản có attention đúng bằng phần trọng số của \texttt{BahdanauAttention}
($W_h, W_f, v$), phần còn lại (embedding tie, LSTMCell, hai lớp khởi tạo
$h_0,c_0$) giữ nguyên. Transformer có xấp xỉ \textbf{2,3×} số tham số của
LSTM+Bahdanau dù cùng $d_{\text{model}}=512$: chi phí nằm ở ba khối lặp lại,
mỗi khối mang hai tầng \texttt{nn.MultiheadAttention} (self + cross, mỗi tầng
có $4$ ma trận chiếu $Q,K,V,O$ kích thước $512\times512$) và một FFN
$512\to2048\to512$ --- so với một \texttt{LSTMCell} duy nhất của decoder kia.

\section{Chi tiết cài đặt hai decoder}
\label{sec:decoder-detail}

Mục này đi sâu vào từng phép biến đổi tensor bên trong hai decoder, khớp trực
tiếp với mã nguồn \texttt{final/models/lstm\_decoder.py} và
\texttt{final/models/transformer\_decoder.py}. Ký hiệu: $B$ là batch size, $R=49$
số vùng ảnh, $D=512=d_{\text{model}}$, $T$ độ dài chuỗi caption tại thời điểm
đang xử lý.

\subsection{LSTM + Bahdanau: từng bước một}

Đầu vào cho decoder là \texttt{feats\_proj} $\in \mathbb{R}^{B\times R\times D}$
(đã qua \texttt{ImageProjection}). Trạng thái ẩn và trạng thái cell của
\texttt{LSTMCell} \emph{không} khởi tạo bằng không mà bằng trung bình các vùng
ảnh, qua hai lớp tuyến tính riêng:
\begin{equation}
	\bar f = \frac{1}{R}\sum_{i=1}^{R} f_i \in \mathbb{R}^{B\times D}, \qquad
	h_0 = \tanh(W_{c\to h}\,\bar f), \qquad
	c_0 = \tanh(W_{c\to c}\,\bar f),
\end{equation}
với $W_{c\to h}, W_{c\to c}: \mathbb{R}^{D}\to\mathbb{R}^{D}$ (\texttt{init\_h},
\texttt{init\_c} trong code) --- cách "gieo mầm" trạng thái ẩn bằng chính nội
dung ảnh, thay vì để mô hình tự học từ vector không, giúp bước sinh từ đầu tiên
đã có ngữ cảnh thị giác.

Tại bước $t$, attention Bahdanau (công thức~\ref{eq:bahdanau}) dùng $h_{t-1}
\in \mathbb{R}^{B\times D}$ và toàn bộ $R$ vùng để tính vector ngữ cảnh $c_t \in
\mathbb{R}^{B\times D}$ và trọng số $\alpha_t \in \mathbb{R}^{B\times R}$ (tổng
theo chiều $R$ bằng 1 --- bất biến được khóa bằng unit test). Đầu vào của
\texttt{LSTMCell} là phép ghép nối (concat), không phải cộng:
\begin{equation}
	x_t = \big[\,e_{t-1} \;;\; c_t\,\big] \in \mathbb{R}^{B\times 2D}, \qquad
	(h_t, c^{\text{cell}}_t) = \text{LSTMCell}(x_t, (h_{t-1}, c^{\text{cell}}_{t-1})),
\end{equation}
với $e_{t-1} \in \mathbb{R}^{B\times D}$ là embedding của từ thật ở bước trước
(teacher forcing). Đây là lý do \texttt{LSTMCell} được khai báo với chiều đầu
vào $2D=1024$ dù hidden size chỉ $D=512$. Logits bước $t$ là
$\text{fc}(\text{dropout}(h_t)) \in \mathbb{R}^{B\times V}$ với $V=2.539$. Biến
thể tắt attention (ablation) thay $c_t$ bằng $\bar f$ cố định và $\alpha_t$
bằng phân phối đều $1/R$ tại mọi bước --- cùng luồng tensor, chỉ khác nguồn
ngữ cảnh.

\subsection{Transformer decoder: từng sublayer trong một khối pre-norm}

Chuỗi từ vào được cộng embedding token với \emph{positional embedding học
được} (không phải sin/cos cố định như bản gốc \autocite{vaswani2017attention}
--- hợp lý hơn khi $T$ luôn $\leq 22$ và có đủ dữ liệu để học bảng vị trí):
\begin{equation}
	x^{(0)} = \text{dropout}\big(E[\text{cap\_in}] + P[0..T{-}1]\big) \in
	\mathbb{R}^{B\times T\times D},
\end{equation}
với $E \in \mathbb{R}^{V\times D}$ là bảng embedding (tie với tầng ra, xem
dưới) và $P \in \mathbb{R}^{22\times D}$ là bảng vị trí (\texttt{max\_len=22}
$=$ 20 từ $+$ \texttt{<bos>} $+$ \texttt{<eos>}).

Mỗi khối (3 khối, lặp lại độc lập tham số) áp dụng ba sublayer theo kiểu
\emph{pre-norm} --- LayerNorm đứng \emph{trước} phép biến đổi, residual cộng
\emph{sau}, khác với bản Transformer gốc (post-norm):
\begin{align}
	x' &= x + \text{dropout}\big(\text{MHA}_{\text{self}}(\text{LN}_1(x),\,
	\text{LN}_1(x),\, \text{LN}_1(x);\ \text{mask}=M)\big), \\
	x'' &= x' + \text{dropout}\big(\text{MHA}_{\text{cross}}(\text{LN}_2(x'),\,
	\text{mem},\, \text{mem})\big), \\
	x''' &= x'' + \text{dropout}\big(\text{FFN}(\text{LN}_3(x''))\big),
\end{align}
trong đó $\text{mem} = \text{feats\_proj} \in \mathbb{R}^{B\times R\times D}$
không đổi qua các khối (đầu ra của \texttt{ImageProjection}, không phải đầu ra
khối trước). Self-attention có $Q,K,V \in \mathbb{R}^{B\times T\times D}$, chia
8 đầu $\Rightarrow$ mỗi đầu $d_k = D/8 = 64$; cross-attention có $Q \in
\mathbb{R}^{B\times T\times D}$ (từ chuỗi chữ) nhưng $K, V \in
\mathbb{R}^{B\times R\times D}$ (từ 49 vùng ảnh) --- đúng vai trò khối
\emph{fusion} đa phương thức. Trọng số cross-attention trung bình 8 đầu, hình
$(B, T, R)$, được giữ lại từ khối cuối cùng để vẽ bản đồ nhiệt.

Mặt nạ nhân quả $M \in \{0,1\}^{T\times T}$ là tam giác trên nghiêm ngặt
(\texttt{torch.triu(\dots, diagonal=1)}), giá trị \texttt{True} tại $(i,j)$
với $j>i$ nghĩa là vị trí $i$ \emph{không được phép} nhìn vị trí $j$ --- chặn
nhìn tương lai trong khi vẫn xử lý toàn bộ $T$ vị trí song song (khác LSTM
phải lặp tuần tự $T$ bước). FFN là $D \to 2048 \to D$ với GELU ở giữa. Vì
kiến trúc pre-norm không chuẩn hoá đầu ra khối cuối, cần một
\textbf{LayerNorm cuối} ($\text{LN}_{\text{final}}$) trước khi chiếu ra logits
--- thiếu tầng này, phân phối $x^{(3)}$ đưa thẳng vào \texttt{fc} sẽ có
phương sai trôi dạt tuỳ theo độ sâu, không so sánh được giữa các vị trí:
\begin{equation}
	\text{logits} = \text{LN}_{\text{final}}(x^{(3)})\, E^\top \in
	\mathbb{R}^{B\times T\times V}.
\end{equation}

\subsection{Weight tying và hệ quả lên khởi tạo}

Cả hai decoder gán \texttt{fc.weight = embedding.weight}: tầng chiếu ra logits
dùng lại đúng ma trận embedding thay vì một ma trận $D\times V$ độc lập
\autocite{press2017using}. Lợi ích: giảm khoảng 1,5 triệu tham số ở vocab
2.539 từ (bằng đúng $D \times V = 512 \times 2.539$), và ràng buộc không gian
biểu diễn của từ ở đầu vào phải nhất quán với không gian dùng để xếp hạng từ ở
đầu ra --- quan trọng khi dữ liệu huấn luyện nhỏ (30.000 caption). Cái giá
phải trả là ma trận embedding không còn "chỉ là embedding": nó đồng thời đóng
vai trò ma trận chiếu đầu ra, nên thang khởi tạo của nó ảnh hưởng trực tiếp
đến phương sai của logits ngay tại bước đầu tiên. Mục~\ref{sec:embedding-init}
(chương~\ref{chap:experiments}) trình bày chi tiết một lỗi thực tế gặp phải vì
điều này và cách phát hiện/sửa nó.

\section{Cấu hình huấn luyện}

\begin{table}[H]
	\centering
	\caption{Siêu tham số huấn luyện (chung cho mọi run trừ khi ghi khác).}
	\label{tab:hyperparams}
	\begin{tabular}{ll}
		\hline
		Optimizer & AdamW, weight decay $0{,}01$ \\
		Batch size & 128 \\
		Learning rate & LSTM: $3\times10^{-4}$; Transformer: $10^{-4}$, warmup 2.000 bước \\
		Lịch LR & cosine decay theo bước \\
		Label smoothing & 0,1 \\
		Gradient clipping & 5,0 \\
		Số epoch tối đa & 25, early stopping theo val loss (patience 5) \\
		Seed & 42 (cố định cho mọi run) \\
		\hline
	\end{tabular}
\end{table}

Suy luận dùng greedy và beam search $B \in \{3, 5\}$ với chuẩn hóa độ dài.
Giai đoạn SCST (mở rộng) khởi động từ checkpoint SFT tốt nhất, phần thưởng là
CIDEr, baseline là điểm của câu giải mã greedy.

\section{Môi trường thực nghiệm và tổ chức mã nguồn}

Mã nguồn tổ chức thành package \texttt{final/} (PyTorch), mọi lựa chọn thí nghiệm
đi qua một \texttt{Config} dataclass trung tâm --- đổi thí nghiệm là đổi cờ dòng
lệnh, không sửa code (cùng khuôn với đồ án giữa kỳ). Kiểm thử đơn vị (pytest)
khóa các bất biến: mặt nạ nhân quả không nhìn tương lai, trọng số attention có
tổng bằng 1, encode/decode từ điển khứ hồi, greedy/beam dừng đúng ở
\texttt{<eos>}, và (bổ sung sau phát hiện ở
mục~\ref{sec:embedding-init}) thang khởi tạo embedding đủ nhỏ khi tie weight.

\subsection{Hai máy, hai vai trò}

\textbf{Máy phát triển (Mac, chip Apple Silicon, backend MPS)} dùng để viết
code theo TDD và chạy \emph{smoke test}: mỗi lệnh huấn luyện có cờ
\texttt{--smoke} rút tập huấn luyện còn 512 mẫu, 2 epoch --- đủ để xác nhận
toàn bộ pipeline (tải dữ liệu $\to$ build vocab $\to$ trích đặc trưng $\to$
train $\to$ evaluate $\to$ visualize) chạy hết không lỗi trong vài phút, trước
khi tốn thời gian trên máy có GPU. \texttt{pick\_device()} tự chọn CUDA
$\to$ MPS $\to$ CPU theo thứ tự ưu tiên nên cùng một lệnh chạy được trên cả
hai máy.

\textbf{Máy huấn luyện thật (``trainbox'')} là máy Windows chạy WSL2 (Ubuntu),
GPU NVIDIA RTX PRO 4000 Blackwell 24\,GB, Python 3.12 quản lý bằng \texttt{uv}
trong virtualenv riêng (\texttt{\textasciitilde/work/.venv}), PyTorch~2.11 với
CUDA~12.8 (\texttt{cu128}) --- kiến trúc Blackwell \emph{bắt buộc} cu128 trở
lên, bản cu121 không nhận diện được GPU này. Truy cập qua SSH; tiến trình
huấn luyện chạy trong \texttt{tmux} để không mất phiên khi mất kết nối mạng,
kéo checkpoint và output về máy dev bằng \texttt{rsync}.

\subsection{Quy trình phát triển: TDD theo 13 task, review độc lập}

Toàn bộ package được lập kế hoạch trước thành 13 task tuần tự (từ scaffold
\texttt{Config} đến SCST), mỗi task viết theo chu trình \emph{test trước, code
sau}: định nghĩa hợp đồng đầu vào/ra của module, viết unit test cho hợp đồng
đó, thấy test đỏ, rồi mới cài đặt tới khi test xanh. Tổng cộng 45 unit test
được viết qua quá trình này (phụ lục~\ref{app:tests} liệt kê đầy đủ theo nhóm),
phủ từ tokenize/
vocab, hình dạng tensor của từng decoder, mặt nạ nhân quả, đến beam search
dừng đúng \texttt{<eos>}. Mỗi task, sau khi code xanh, còn qua một bước
\emph{review độc lập} (một phiên riêng đọc lại diff so với hợp đồng đã định
nghĩa, không dùng lại ngữ cảnh của phiên viết code) trước khi coi là xong ---
mục tiêu là bắt các lỗi kiểu ``code chạy nhưng sai hợp đồng'' mà bản thân tác
giả code dễ bỏ sót vì đã quen với logic vừa viết.

\subsection{\texttt{run\_all.sh}: idempotent, chống đứt phiên}

Pipeline đầy đủ (tải dữ liệu $\to$ vocab $\to$ trích đặc trưng $\to$ ba run
huấn luyện chính $\to$ evaluate $\to$ visualize) được gói trong một script duy
nhất, chạy được cả ở chế độ \texttt{--smoke} lẫn chạy thật. Script không giả
định chạy trọn vẹn một lần: trước mỗi bước huấn luyện, nó kiểm tra checkpoint
tương ứng (\texttt{final/checkpoints/<run>.pt}) đã tồn tại chưa --- nếu có thì
bỏ qua bước đó. Nhờ vậy, nếu phiên SSH tới trainbox bị đứt giữa chừng (mất
mạng, hết giờ), chạy lại đúng lệnh cũ sẽ tiếp tục từ run chưa xong thay vì làm
lại từ đầu; đây là lý do thiết kế lưu checkpoint theo tên run cố định
(\texttt{lstm.pt}, \texttt{transformer.pt}, \dots) thay vì theo timestamp.
